% =====================================================================
%  CARNET D'ENTRAÎNEMENT — Fiche 12 (Chimie organique)
%  THÈME 2 — Reconnaissance et classification des groupements fonctionnels
%  NIVEAU MOYEN — CORRIGÉ
% =====================================================================
\documentclass[11pt,a4paper]{article}
\usepackage[utf8]{inputenc}
\usepackage[T1]{fontenc}
\usepackage{lmodern}
\usepackage[french]{babel}
\usepackage{amsmath,amssymb,mathtools}
\providecommand{\degree}{^\circ}
\usepackage{icomma}
\usepackage{siunitx}
\sisetup{output-decimal-marker={,},per-mode=symbol}
\usepackage[version=4]{mhchem}
\usepackage{chemfig}
\setchemfig{atom sep=2em,bond length=0.9em,cram width=2.5pt}
\usepackage{xcolor}
\usepackage{array}
\usepackage{booktabs}
\usepackage{tabularx}
\usepackage[most]{tcolorbox}
\usepackage{enumitem}
\usepackage{tikz}
\usetikzlibrary{arrows.meta,positioning,calc}
\usepackage{fancyhdr}
\usepackage{titlesec}
\usepackage[hidelinks]{hyperref}
\usepackage[a4paper,margin=2.1cm,headheight=26pt,headsep=9pt]{geometry}

\definecolor{primary}{RGB}{150,98,162}
\definecolor{primarydark}{RGB}{92,52,110}
\definecolor{excol}{RGB}{0,135,90}
\definecolor{attncol}{RGB}{200,40,55}
\definecolor{methcol}{RGB}{90,90,100}

\newcommand{\runtitle}{Thème 2 — Moyen — Corrigé}
\pagestyle{fancy}\fancyhf{}
\fancyhead[L]{\small\textcolor{primarydark}{\textbf{Fiche 12} — Chimie organique}}
\fancyhead[R]{\small\textcolor{primarydark}{\runtitle}}
\fancyfoot[C]{\small\thepage}
\renewcommand{\headrulewidth}{0.8pt}
\renewcommand{\headrule}{\hbox to\headwidth{\color{primary}\leaders\hrule height \headrulewidth\hfill}}
\setlength{\parindent}{0pt}\setlength{\parskip}{3pt}

\tcbset{boxbase/.style={enhanced,breakable,boxrule=0.9pt,arc=2.5pt,
   left=3mm,right=3mm,top=2mm,bottom=2mm,fonttitle=\bfseries,coltitle=white,
   toptitle=0.5mm,bottomtitle=0.5mm}}
\newtcolorbox[auto counter]{corr}[1]{boxbase,colback=excol!5,colframe=excol,
   title={\textsc{Corrigé}~\thetcbcounter\ \textmd{--- \itshape #1}}}
\newcommand{\rep}[1]{\textcolor{primarydark}{\textbf{#1}}}

\newcommand{\banner}[2]{%
\begin{tcolorbox}[enhanced,colback=excol,colframe=excol,arc=5pt,boxrule=0pt,
   left=5mm,right=5mm,top=2.5mm,bottom=2.5mm,coltext=white]
{\large\bfseries #1}\\[1pt]{\small #2}
\end{tcolorbox}\vspace{2pt}}

\begin{document}

\banner{Thème 2 — Reconnaissance des groupements fonctionnels}{Niveau \textbf{Moyen} — Corrigé détaillé \quad$\vert$\quad Fiche 12, Chimie organique}

% ---------------------------------------------------------------------
\begin{corr}{recenser toutes les fonctions}
On parcourt $\ce{HO-CH2-CO-CH2-COOH}$ :
\begin{itemize}[leftmargin=1.6em,itemsep=1pt]
  \item $\ce{HO-CH2-}$ : O--H sur un C saturé (lié à 1 autre C) $\Rightarrow$ \rep{alcool primaire} ;
  \item le \textbf{premier} carbonyle est encadré par deux carbones $\Rightarrow$ \rep{cétone} ;
  \item le \textbf{second} carbonyle porte un O--H $\Rightarrow$ \rep{acide carboxylique}.
\end{itemize}
Bilan : \rep{3 fonctions} — un alcool (primaire), une cétone, un acide carboxylique. Les \textbf{deux} $\ce{C=O}$ appartiennent à des fonctions différentes selon leur voisinage (2 C $\to$ cétone ; O--H $\to$ acide).
\end{corr}

% ---------------------------------------------------------------------
\begin{corr}{compter et classer les alcools}
\textbf{a)} La molécule porte \rep{trois} groupes $\ce{-OH}$, donc \rep{3 fonctions alcool}.\\
\textbf{b)} On examine le carbone porteur de chaque $\ce{-OH}$ :
\begin{itemize}[leftmargin=1.6em,itemsep=1pt]
  \item $\ce{HO-CH2-}$ : le C est lié à 1 autre C $\Rightarrow$ \rep{alcool primaire} ;
  \item le C central $\ce{C(OH)(CH3)}$ : lié à 3 autres C (le $\ce{CH2}$, le $\ce{CH3}$, le $\ce{CH}$) $\Rightarrow$ \rep{alcool tertiaire} ;
  \item le $\ce{CH(OH)}$ : lié à 2 autres C (le C central et le $\ce{CH3}$ terminal) $\Rightarrow$ \rep{alcool secondaire}.
\end{itemize}
Bilan : \rep{1 primaire, 1 secondaire, 1 tertiaire}. On compte \emph{et} on classe.
\end{corr}

% ---------------------------------------------------------------------
\begin{corr}{discriminer les carbonyles}
\textbf{a)} \rep{Deux} groupes carbonyle $\ce{C=O}$.\\
\textbf{b)} On regarde le voisinage de chacun :
\begin{itemize}[leftmargin=1.6em,itemsep=1pt]
  \item $\ce{CH3-CO-CH2-}$ : carbonyle entre deux carbones $\Rightarrow$ \rep{cétone} ;
  \item $\ce{-CHO}$ (bout de chaîne, porte un H) $\Rightarrow$ \rep{aldéhyde}.
\end{itemize}
\textbf{c)} Fonctions totales : \rep{une cétone, une double liaison $\ce{C=C}$ (alcène), un aldéhyde}. \rep{Aucune fonction acide carboxylique} (aucun $\ce{C=O}$ ne porte de O--H).
\end{corr}

% ---------------------------------------------------------------------
\begin{corr}{azote : amine ou amide, et de quelle classe ?}
\textbf{a)} et \textbf{b)} — on parcourt $\ce{H2N-CH2-CH2-CO-NH-CH2-CH2-N(CH3)2}$ :
\begin{itemize}[leftmargin=1.6em,itemsep=1pt]
  \item $\ce{H2N-CH2-}$ : azote lié à un C saturé, sans carbonyle $\Rightarrow$ \rep{amine primaire} ;
  \item $\ce{-CO-NH-}$ : azote accolé à un carbonyle $\Rightarrow$ \rep{amide} (ce n'est \emph{pas} une amine, malgré la liaison $\ce{C-N}$) ;
  \item $\ce{-N(CH3)2}$ : azote lié à 3 carbones (le $\ce{CH2}$ et deux $\ce{CH3}$), sans carbonyle $\Rightarrow$ \rep{amine tertiaire}.
\end{itemize}
Bilan : \rep{une amine primaire, une amide, une amine tertiaire}. Le piège classique — appeler l'amide « amine secondaire » — est évité car un $\ce{C=O}$ est accolé à cet azote.
\end{corr}

% ---------------------------------------------------------------------
\begin{corr}{effet d'une réduction sur le décompte des fonctions}
\textbf{a)} Fonctions oxygénées de $(\ce{CH3})_2\ce{C(OH)-CH2-CO-CH2-CH(OH)-CH3}$ :
\begin{itemize}[leftmargin=1.6em,itemsep=1pt]
  \item $(\ce{CH3})_2\ce{C(OH)-}$ : C--OH lié à 3 autres C $\Rightarrow$ \rep{alcool tertiaire} ;
  \item $\ce{-CO-}$ entre deux $\ce{CH2}$ $\Rightarrow$ \rep{cétone} ;
  \item $\ce{-CH(OH)-CH3}$ : C--OH lié à 2 autres C $\Rightarrow$ \rep{alcool secondaire}.
\end{itemize}
Soit \rep{2 alcools (1 tertiaire, 1 secondaire) + 1 cétone}.\\
\textbf{b)} La réduction transforme $\ce{-CO-}$ en $\ce{-CH(OH)-}$. Ce nouveau carbone porteur du $\ce{-OH}$ est encadré par deux $\ce{CH2}$, soit 2 autres C $\Rightarrow$ nouvel \rep{alcool secondaire}. La cétone disparaît, remplacée par un alcool. Total : \rep{3 fonctions alcool} (1 tertiaire + 2 secondaires) et \rep{plus aucune cétone}.
\end{corr}

\end{document}
